\documentclass[11pt,a4paper]{article}
\usepackage{fontspec}
\setmainfont{Noto Serif}
\setsansfont{Noto Sans}
\usepackage{babel}
\babelprovide[main,import]{russian}
\usepackage{amsmath,amssymb,mathtools,booktabs,geometry,fancyhdr,microtype,xcolor,hyperref}
\geometry{left=21mm,right=21mm,top=21mm,bottom=21mm}
\definecolor{ink}{HTML}{173D57}
\hypersetup{colorlinks=true,urlcolor=ink,linkcolor=ink,pdftitle={BGS-F. Перестановки участков полного обхода}}
\pagestyle{fancy}\fancyhf{}\lhead{\small BGS-F}\rhead{\small Перестановки участков полного обхода}\cfoot{\thepage}
\setlength{\headheight}{15pt}\setlength{\parindent}{1em}\setlength{\parskip}{2.5pt}
\setlength{\emergencystretch}{3em}\tolerance=3000
\newcommand{\Tr}{\operatorname{Tr}}
\newcommand{\EE}{\mathbb E}
\newcommand{\R}{\mathbb R}
\newcommand{\Z}{\mathbb Z}
\newcommand{\C}{\mathbb C}
\newcommand{\si}{\mathrm{Sine}_2}
\newcommand{\eps}{\varepsilon}
\newcommand{\sinc}{\operatorname{sinc}}
\newcommand{\sect}[1]{\section*{\sffamily\color{ink}#1}}
\newcommand{\subt}[1]{\par\medskip\noindent\textbf{#1.}\quad}
\newcommand{\T}{\mathbb T}
\newcommand{\dist}{\operatorname{dist}}
\newcommand{\Ad}{\operatorname{Ad}}
\begin{document}
\noindent{\small\sffamily\color{ink} МАТЕМАТИЧЕСКАЯ ФИЗИКА\hfill 20 НОЯБРЯ 2026}
\begin{center}{\LARGE\sffamily Перестановки участков полного обхода}\\[4pt]
{\large Знаковое суммирование по растущему числу участков}
\end{center}
Для исходной карты $U_N=e^{i10^{-80}NA_N}M_N$ рассмотрим $p_m=\Tr U_N^m$. В разложении $|p_m|^2$ пара замкнутых слов $(x,y)$ длины $m$ имеет вес $a(x)\overline{a(y)}$, где $a(x)=\prod_{j=0}^{m-1}U_{x_{j+1},x_j}$. Слова укоренены в позиции 0. В обоих словах зададим $r$ коротких участков, разделённых свободными промежутками. Потребуем, чтобы участки во втором слове были перестановкой участков первого с сохранением корневого участка. Для подходящих пар определяющим условием служат только эти списки вершин. Сумма по классу с учётом каждой пары ровно один раз равна
\[
  (1+O(N^{-5}))\,\frac{N}{r(r+1)}
       \left[\prod_{j=0}^{r}\left(1+\frac jN\right)
                    -\prod_{j=0}^{r}\left(1-\frac jN\right)\right].
                                                                     \tag{1}
\]
Длина короткого участка $\ell=\lceil8\log_2 N\rceil$. На полосе $m\asymp N$ оценка допускает $r$ порядка $\sqrt{N\log N}$. При таком выборе класс покрывает перестановки по всему обходу, при этом суммарная длина отмеченных участков $r\ell=o(m)$. Каждый свободный промежуток также растёт с $N$.

При $r(r+1)/(2N)=\log N+\log\log N+\log(2c)+o(1)$, вклад класса равен $cN+o(N)$ для любого фиксированного $c>0$. Дополнение внутри всех пар с общим корневым участком даёт $-cN+o(N)$. Для второго класса знак установлен вычислением, без предположения о статистике фаз. Выбор $r$ в зависимости от $m/N$ позволяет получить в этой отрицательной части профиль $-(m-N)_+$ с равномерной ошибкой $o(N)$. Положительная часть имеет противоположный профиль. Оба класса вместе дают вклад $o(1)$ после нормировки и усреднения в тесте (P).

Доказательство использует смешивание интенсивностей на коротких совпавших участках и унитарность на свободных промежутках. Вклад заданной перестановки определяется числом циклов произведения двух циклических порядков. Формула суммирования этих чисел известна в комбинаторике симметрических групп [1]. Ниже выведены её связь со знаковой траекторной суммой и равномерная оценка ошибок, в том числе при совпадении нескольких участков.

\textbf{Статус.} Вычислены макроскопические по величине вклады дальних перестановок при $m\asymp N$. Полный парный тест и сходимость к $\si$ остаются открытыми.
\newpage
\sect{1. Геометрия класса и унитарное стягивание}
Работаем при $m\le2N$ и $r(\ell+1)\le m$. Положим $g=\lfloor m/r\rfloor-\ell\ge1$, $q=m-r(\ell+g)$, $g_0=g+q$, $g_i=g$ для $i>0$. В первом слове в позиции 0 начинается участок 0 длины $\ell$ рёбер. После участка $i$ следует свободный промежуток длины $g_i$ рёбер и следующий участок. Порядок меток даёт цикл $\kappa=(0\,1\,\dots\,r-1)$. Во втором слове участки стоят на тех же временных позициях. Корневой участок совпадает, а остальные участки требуем сопоставить биективно с участками первого слова. Порядок второго обхода после разметки даёт некоторый $r$-цикл $\sigma$. После метки 0 во втором слове также идёт промежуток $g+q$; после всех других --- $g$. Поэтому при каждом сопоставлении длина промежутка после одной и той же метки одинакова в обоих словах.

Пусть $B_{ab}=|U_{ab}|^2$, $J_{ab}=1/N$. Из (7.8) постановки и предыдущей части используем
\[
       B^\ell=J+E,\qquad\|E\|\le2^{-\ell},\qquad
       \max_{a,b}B_{ab}\le K/N,\quad K\ge1.                           \tag{2}
\]
Через $a_i,b_i$ обозначим начальную и конечную вершины общего участка $i$. При заранее заданном сопоставлении, без запрета на повторы, суммирование по внутренним вершинам участков и промежутков даёт
\[
 \sum_{a_0,\ldots,a_{r-1}}\ \prod_{i=0}^{r-1}
          \sum_{b_i}
          (B^\ell)_{b_i a_i}
          (U^{g_i})_{a_{\kappa(i)},b_i}
           \overline{(U^{g_i})_{a_{\sigma(i)},b_i}} .                 \tag{3}
\]
Множитель $J$ в позиции $i$ даёт $\mathbf1_{a_{\kappa(i)}=a_{\sigma(i)}}/N$ по унитарности $U^{g_i}$. Все равенства сшивают вершины по циклам перестановки $\kappa\sigma^{-1}$. Пусть $c(\pi)$ означает число циклов, включая одноэлементные. Главный член (3) равен
\[
                             N^{c(\kappa\sigma^{-1})-r}.
                                                                     \tag{4}
\]

Если в $s$ позициях взять $E$, то исчезает $s$ равенств и добавляется не более $s$ свободных меток. В каждой такой позиции модуль суммы по $b_i$ не превосходит $2^{-\ell}$ по Коши и нормировке строк $U^{g_i}$. Модуль ошибки в (3) поэтому ограничен главным членом, умноженным на $(1+\beta)^r-1$, где $\beta=N^2 2^{-\ell}$. Эта оценка относительна отдельно для каждого $\sigma$. Суммирование по $(r-1)!$ порядкам сохраняет ту же относительную точность.

Для двух пересекающихся факторов с различными длинами промежутков аналогичная подстановка $J$ стягивает произведение в $U^{g_i^x-g_i^y}$, а не в тождественную матрицу. Вдоль каждой компоненты появляются следы степеней $U$. Согласование длин после меток в данном классе устраняет необходимость в оценках этих следов на больших временах.
\newpage
\sect{2. Повторные участки и единственность сопоставления}
Совпадением двух участков называем равенство полных списков из $\ell+1$ вершин. Обозначим список в метке $i$ через $w_i$. Если все $w_i$ различны, то из пары $(x,y)$ однозначно восстанавливается $\sigma$.

Зафиксируем разбиение $\mathcal P$ множества меток, потребовав равенство $w_i$ внутри каждой его части. Для части размера $k\ge2$ общий множитель от участков, с учётом их совпадающих концов, равен $(B^{\circ k})^\ell_{ba}$, где $B^{\circ k}$ означает покомпонентную степень. По (2) и стохастичности
\[
     (B^{\circ k})^\ell_{ba}
        \le (K/N)^{\ell(k-1)} B^\ell_{ba}
        \le (K/N)^{\ell(k-1)} K/N .                                 \tag{5}
\]
Сумма по общему $b$ произведения модулей $2k$ элементов строк унитарных матриц не превосходит 1 по Гёльдеру. Для оценки удаляем $k$ связей от $J$ и учитываем рост числа свободных меток не более чем на $k$. Соответствующий относительный множитель ограничен
\[
     K N^{2k-1}(K/N)^{\ell(k-1)}
                 \le \theta^{k-1},\qquad \theta=K N^3(K/N)^\ell .
\]
Метка начала также отождествляется внутри части. Снятие этого ограничения при оценке модуля может только увеличить верхнюю границу. Связи от остальных частей размера один обрабатываем как в (3). Для всех частей размера не меньше двух относительные множители перемножаются.

Индикатор попарно различных $w_i$ раскладывается по разбиениям на суммы с коэффициентами $\prod_{C\in\mathcal P}(-1)^{|C|-1}(|C|-1)!$. Это включение--исключение можно получить, суммируя знаки перестановок, чьи циклы лежат внутри классов равенства $w_i$: каждый нетривиальный класс даёт ноль. Абсолютные коэффициенты при $s=r-\#\mathcal P$ суммируются в степени $z^s$ по формуле
\[
                         \prod_{j=0}^{r-1}(1+jz).                    \tag{6}
\]
Действительно, каждая часть размера $k$ допускает $(k-1)!$ циклических порядков. Новый элемент либо образует отдельный цикл, либо вставляется за одним из прежних элементов.
Поэтому для класса попарно различных участков (4) сохраняется для каждого $\sigma$ с относительной ошибкой не более
\[
              (1+\beta)^r\prod_{j=0}^{r-1}(1+j\theta)-1=O(N^{-5}).
                                                                     \tag{7}
\]

То же верно для полной суммы по биективно сопоставимым участкам с учётом каждой пары один раз. Для проверки сгруппируем слова по мультимножеству $M$ участков и слову на корневом участке. Если $n_w$ --- кратности среди \emph{некорневых} участков, то один класс имеет знаковую массу $|\sum_x a(x)|^2\ge0$. В сумме по всем сопоставлениям класс имеет ту же массу, умноженную на $\prod_w n_w!$. В группе с повтором, включая повтор корневого слова, её немеченый вклад заключён между нулём и её вкладом с сопоставлениями. Последний ограничен разностью (3) и результата с попарно различными участками, после суммирования по $\sigma$, то есть $O(N^{-5})$ от суммы (4). Так установлена оценка одновременно для попарно различных участков и для класса, допускающего повторы.
\newpage
\sect{3. Сумма по порядкам и её масштаб}
После (7) остаётся комбинаторическая сумма $S_{r,N}=N^{-r}\sum_{\sigma}N^{c(\kappa\sigma^{-1})}$ по всем $r$-циклам. Приведём вывод формулы Стэнли [1] в нужном частном случае. Характер представления $S_r$ перестановками факторов тензорной степени $(\C^N)^{\otimes r}$ равен $N^{c(\pi)}$. По двойственности Шура--Вейля кратность компонент типа $\lambda$ равна $s_\lambda(1^N)$. Сумма по классу $r$-циклов действует в неприводимой компоненте скаляром $(r-1)!\,\chi_\lambda(\kappa)/\dim\lambda$. Характер длинного цикла отличен от нуля только на крюках $\lambda=(r-k,1^k)$, где он равен $(-1)^k$ [1,2]. Значит,
\[
    \sum_{\sigma}N^{c(\kappa\sigma^{-1})}
       =(r-1)!\sum_{k=0}^{r-1}
           \frac{s_{(r-k,1^k)}(1^N)}{\binom{r-1}{k}}
       =\frac1r\sum_{k=0}^{r-1}\prod_{j=-k}^{r-1-k}(N+j).
\]
Второе равенство --- формулы крюков и содержаний. Каждый член произведения есть разность произведений длины $r+1$ с нижними пределами $-k$ и $-k-1$, делённая на $r+1$. Телескопирование даёт (1).

При $\mu=r(r+1)/(2N)$ и $r=o(N^{2/3})$,
\[
             S_{r,N}=(1+o(1))\,\frac{\sinh\mu}{\mu}.                  \tag{8}
\]
В самом деле $\log\prod_{j=0}^r(1\pm j/N)=\pm\mu+O(r^3/N^2)$. При $\mu\to0$ ошибка после вычитания имеет относительный порядок $O(r^3/(N^2\mu))=O(r/N)$, при $\mu$ отделённом от нуля относительная ошибка $o(1)$ непосредственно. Формула (8) применима с относительной ошибкой в области $\mu\asymp\log N$. Из (7),(8) следует утверждение первого листа для произвольного $c>0$.

\subt{Распределение числа связей} Для попарно различных участков можно отдельно суммировать по перестановкам с любым заданным числом $d=r-c(\kappa\sigma^{-1})$. После нормирования главные члены образуют вероятность, а её производящая функция равна $S_{r,N/z}/S_{r,N}$ при $z>0$ (выражение понимается как многочлен по $z$). Ошибка для каждого множества порядков имеет относительную оценку (7). При $\mu\to u>0$ функция стремится к $(\sinh(uz)/(uz))/(\sinh u/u)$. В ростовом режиме $\mu\asymp\log N$ величина $(d-\mu)/\sqrt\mu$ стремится под этой вероятностью к стандартному нормальному закону. Подстановка $z=e^{s/\sqrt\mu}$ в (8) даёт предельную функцию моментов $e^{s^2/2}$ после центрирования. Число $d$ чётно, поскольку $\kappa$ и $\sigma$ имеют один знак. В режиме вклада порядка $N$ суммирование только по ограниченному числу разрывов $d$ упускает основную массу.
\newpage
\sect{4. Прямое сравнение с форм-фактором}
Обозначим через $\mathcal C_{\ell}$ все пары с общим корневым участком и через $\mathcal A_{r,\ell}$ её класс с перестановками из раздела 1. Пары из $\mathcal C_{\ell}$ суммируются точно:
\[
       \sum_{\mathcal C_{\ell}} a(x)\overline{a(y)}
             =\Tr (B^\ell C_{m-\ell}),\qquad
       (C_j)_{ab}=|(U^j)_{ab}|^2,\quad
       |\Tr (B^\ell C_{m-\ell})-1|\le N2^{-\ell}.
\]
Следовательно, сумма на $\mathcal C_\ell\smallsetminus\mathcal A_{r,\ell}$ равна $1-S_{r,N}+O(N^{-5}S_{r,N}+N2^{-\ell})$. При (8) с $\mu=\log N+\log\log N+\log(2c)+o(1)$ оба семейства дают конечный ненулевой вклад в исходной нормировке $N^{-2}\sum_m w(m/N)$. Соответствующие пределы равны $c\int_0^2 w(t)\,dt$ и $-c\int_0^2 w(t)\,dt$. Для $m<r(\ell+1)$ принимаем оба класса пустыми. Число таких времён $o(N)$, а оценки на остальных временах равномерны.

Можно получить именно зависимость от времени, необходимую для поправки к рампе. При $N<m<2N$ возьмём ближайшее целое $r$ к решению
\[
  \frac{r(r+1)}{2N}
      =\log N+\log\log N+\log\!\left(2\max\{m/N-1,1/\log N\}\right).
\]
По (8) равномерно $S_{r,N}=(m-N)+o(N)$, так как для $m/N-1<1/\log N$ обе величины $o(N)$. При $m\le N$ примем классы пустыми. Сумма на $\mathcal C_\ell\smallsetminus\mathcal A_{r,\ell}$ теперь даёт профиль $-(m-N)_+$ с равномерным остатком $o(N)$ на $1\le m<2N$. Взвешенный предел равен $-\delta_*$ из (P$'$). Одновременно $\mathcal A_{r,\ell}$ даёт $+\delta_*$. Удаление циклических пар из любого из этих классов меняет его сумму лишь на $O(1)$ равномерно на $N<m<2N$, как показано ниже. Таким образом, оба вклада относятся к собственно нециклической интерференции.

\subt{Циклические пары внутри корневого класса}
При нулевом сдвиге сумма равна $\Tr B^m=1+O(N2^{-m})$. Ненулевой сдвиг $s$ требует в слове $x$ равенств $x_i=x_{i+s}$ при $0\le i\le\ell$, с индексами по модулю $m$. Граф этих равенств имеет ранг не меньше $(\ell+1)/2$: выбранные ориентированные рёбра имеют входящую и исходящую степени не более одной, и полностью замкнутые циклы имеют длину как минимум два. Вес такого события под $Q(x)=|a(x)|^2$ не более $K^{2(\ell+1)}N^{-(\ell+1)/2}$. Действительно, выписываем циклическую маргиналь на затронутых позициях и применяем $\max (B^j)_{ab}\le K/N$ при $j\ge1$ к переходам. Сумма по всем ненулевым сдвигам исчезает быстрее любой степени $N^{-1}$.

\textbf{Статус (P).} Корневой класс целиком даёт $O(1)$ при каждом времени, а его положительная и отрицательная части могут быть порядка $N$ и больше. Вклад пар, у которых корневые участки различны, вместе с совпадениями при других сдвигах остаётся необходим для полного теста.
\newpage
\sect{5. Любое положение длинного совпавшего участка}
Следующее утверждение разрешает различия по любому числу промежутков с произвольными длинами. При \emph{фиксированном прямом совмещении} включим пару один раз, если существует хотя бы один непрерывный общий участок из $\ell$ рёбер в каком-либо месте обхода. Даже при $\ell$ порядка $\log N$ суммарный вклад равен
\[
                 1+O\!\left(N2^{-m}
                        +N\big[(1+2^{-\ell})^m-1\big]\right).         \tag{9}
\]
В частности при $m\asymp N$, $\ell=\lceil8\log_2 N\rceil$ все различные по циклическому классу слова в этом семействе имеют суммарный вклад $O(N^{-6})$. Для $m$ разных совмещений результат (9) справедлив с подсчётом по совмещениям. Его перенос на объединение требует отдельно оценить пересечения несовпадающих совмещений.

Положим $V=U\otimes\bar U$, $P=\sum_a|a,a\rangle\langle a,a|$, $Q=I-P$. Вектор $e=N^{-1/2}\sum_a|a,a\rangle$ неподвижен для $V$, причём $Pe=e$, $Qe=0$. На ортогональном дополнении $PVP$ сжимает норму не менее чем вдвое по (2) с $\ell=1$. Для нетождественных слов общий участок существует тогда и только тогда, когда есть старт $j$ с несовпадением в $j-1$ и совпадениями во всех позициях $j,\ldots,j+\ell$. Объединяем события стартов по включению--исключению. Несколько стартов, отстоящих по кругу менее чем на $\ell+2$ позиций последовательно, несовместимы. Прочие пересечения представлены следом циклического произведения $V$ и проекторов $P,Q$. У каждого из $h$ стартов есть отдельная последовательность $(PVP)^\ell$ между проекторами $P$. Фактор $Q$ убирает прямую $e$, а норма всего произведения на её дополнении не превосходит $2^{-h\ell}$. Его ранг не больше $N$, поскольку присутствует $P$. Сумма абсолютных оценок всех непустых пересечений не больше $N[(1+2^{-\ell})^m-1]$. Тождественные пары дают $\Tr B^m=1+O(N2^{-m})$. Для остальных циклических копий из этого класса используем доказательство раздела 4 и умножаем оценку на $m$ за возможные положения общего участка.

\textbf{Связь с предыдущими оценками.} В предыдущей части контролировались пары с высокой долей совпавших позиций после совмещения или с общим участком длиной больше половины обхода. При фиксированном совмещении (9) заменяет это ограничение логарифмической длиной и допускает произвольное число остальных участков. Формула (1) требует совпадения $r\ell$ рёбер в заданных местах после перестановки участков. Согласование промежутков обеспечивает количественное суммирование по всем перестановкам.

\textbf{Источники.} [1] R. P. Stanley, \emph{Two enumerative results on cycles of permutations}, The Ramanujan Journal 21 (2010). \href{https://math.mit.edu/~rstan/papers/cycles.pdf}{Авторская версия}. [2] W. Fulton, J. Harris, \emph{Representation Theory. A First Course}, GTM 129. \href{https://doi.org/10.1007/978-1-4612-0979-9}{Издание Springer}. Точные матричные и комбинаторные проверки приведены в архиве.
\newpage
\sect{6. Предельное смешивание при всех фиксированных временах}
Из (1)--(9) не получается универсальный переход к полному парному следу. Это можно установить конструктивно на тех же пеллевых размерностях и с той же линейной матрицей $S$. Существует последовательность унитарных матриц $V_N$ с простым спектром, удовлетворяющая следующим трём условиям одновременно:
\[
 \begin{gathered}
       |(V_N)_{ab}|^2=1/N,\\
       \sup_{1\le j<k}\big\|(|(V_N^j)_{ab}|^2)_{ab}-J\big\|\longrightarrow0,\\
       \Delta_{2,N}^{(V)}(b_*)
          \sim \frac{N}{k}\int_0^2 w(t)\,dt\longrightarrow+\infty
                  \quad\text{при }N=h_k,\quad k\longrightarrow\infty
                    \text{ по нечётным простым}.                     \tag{10}
 \end{gathered}
\]
Локальный точечный процесс в этом примере сходится к пустому. Матрицы $V_N$ получены малым диагональным возмущением линейной карты. В основной задаче $U_N$ строится из фиксированного нелинейного кика.

Пусть $h_k+v_k\sqrt2=(1+\sqrt2)^k$, как в предыдущих частях. При нечётном $k$ число $N=h_k$ нечётно и $S^k=I\bmod N$. Если простое $d$ делит $N$, то в кольце $\mathbb F_d[\sqrt2]=\mathbb F_d[X]/(X^2-2)$ элемент $v=3+2\sqrt2$ удовлетворяет $v^k=1$ и $v\ne1$. В самом деле, квадрат $(1+\sqrt2)^k$ равен $1$ в этом кольце по $h_k^2-2v_k^2=-1$. Порядок $v$ равен простому $k$. Для $k\nmid j$ элемент $v^j-1$ обратим: в случае поля это сразу следует из порядка, в случае произведения двух полей компоненты $v$ взаимно обратны по норме $1$, и ни одна из них не может дать единицу в степени $j$. Таким образом $S^j-I$ обратима по модулю $N$ при $k\nmid j$. Перестановка вейлевых мод сопряжением $M_N^j$ даёт
\[
       |\Tr M_N^j|^2=
           \begin{cases}N^2,&k\mid j,\\1,&k\nmid j.\end{cases}        \tag{11}
\]
Второе значение есть ровно один неподвижный вейлев оператор, тождественный.

При $1\le j<k$ внедиагональный элемент $(S^j)_{12}=4h_jv_j$ обратим по модулю $N$. Иначе по некоторому $d\mid N$ элемент $v^j$ был бы скаляром, поскольку его $\sqrt2$-коэффициент обращается в нуль. Его норма $1$ заставила бы этот скаляр равняться $\pm1$, что несовместимо с нечётным простым порядком $k$. Для каждого ненулевого диагонального характера $T_N(0,b)$ сопряжение $M_N^j$ даёт вейлев оператор с ненулевым сдвигом и нулевой диагональю. Следовательно, матрица интенсивностей $|(M_N^j)_{ab}|^2$ равна $J$ для всех $1\le j<k$.
\newpage
\sect{7. Устранение кратности и полный предел примера}
Возьмём $V_N=D_NM_N$, где $D_N$ диагональна и унитарна. Утверждаем, что можно выбрать $\|D_N-I\|\le N^{-3}$ с простым спектром $V_N$. Для каждого собственного подпространства $E$ матрицы $M_N$ с проектором $P_E$ векторы $P_E e_j$ порождают $E$. Значит, в линейном семействе $P_E\operatorname{diag}(d_j)P_E$ с вещественными $d_j$ есть матрица с простым спектром на $E$. Для явного аргумента выберем базис из $\dim E$ таких векторов и последовательно прибавим их положительные внешние произведения, каждый раз с достаточно малым коэффициентом. Ранг повышается на один, прежние положительные собственные значения остаются попарно различными, а новое положительное значение отделено от них. После последнего шага все значения просты. Дискриминант семейства для каждого $E$ тем самым является ненулевым вещественным многочленом. Можно выбрать один $d$ вне объединения конечного числа его нулевых множеств для всех $E$.

Рассмотрим $D(t)=\operatorname{diag}(e^{itd_j})$. Для собственного значения $\lambda$ матрицы $M_N$ первые производные собственных значений $D(t)M_N$ при $t=0$ равны собственным значениям $i\lambda P_E\operatorname{diag}(d_j)P_E$ на $E$. Это следует из первого порядка редукции к спектральному подпространству через контурный проектор. У разных исходных значений есть конечный положительный зазор, а внутри каждого подпространства производные различны. Поэтому при всех достаточно малых положительных $t$ спектр прост. Выбор $t$ с дополнительным ограничением $\|D(t)-I\|\le N^{-3}$ доказывает утверждение.

Левый диагональный множитель сохраняет все модули матричных элементов. Кроме того $\|V_N^m-M_N^m\|\le mN^{-3}$ для любого $m$ по телескопической сумме. Равномерно при $1\le j<k$ это и равенство $|(M_N^j)_{ab}|^2=1/N$ дают вторую строку (10): норма ограничена величиной $2kN^{-5/2}+k^2N^{-5}$ по оценке максимального элемента, умноженной на $N$. При $m\le2N$ след отличается от чистого не более чем на $2/N$. Тогда (11) даёт
\[
  \Delta_{2,N}^{(V)}(b_*)
       =\sum_{\substack{1\le m<2N\\ k\mid m}} w(m/N)+O(1)
       =\frac Nk\left(\int_0^2 w(t)\,dt+o(1)\right).
\]
Вычитание эталона $\min(m,N)$ после деления на $N^2$ ограничено, а ошибка замены следов тоже ограничена. Последний переход есть сумма Римана с шагом $k/N\to0$.

Наконец, спектр $M_N$ лежит на повёрнутой сетке из $k$ узлов. По нормальной резольвенте все фазы $V_N$ лежат в $O(N^{-3})$-окрестности сетки. При равномерном начале развёртки вероятность появления точки в любом фиксированном окне длины $2R$ не больше $k(2R/N+O(N^{-3}))\to0$. Это доказывает слабый предел. Конструкция удовлетворяет точному смешиванию (7.8), сходимости интенсивностей при каждом фиксированном времени и простоте спектра. Для парного теста исходной карты нужно дополнительно использовать фазы её конкретного кика на растущих временах.
\end{document}
